\section{Concluding the Project}

This week concludes the development of my EPQ journey.

My final goal for this week was to produce a video 'walk through' of my completed NeuroNote Artefact which will enable me to demonstrate the features and functionality of my platform.

Thanks to my use of careful planning and template creation during my video presentation production, I was able to re-use many of the same assets and techniques I had already refined such as my gradient background, looping background music and video editing techniques. This meant I was able to focus more on the content of my video and ensuring I clearly demonstrated the features of my platform. 

Unlike my previous presentation I opted not to use a script for this video which enabled me to more naturally demonstrate the features of my platform from a users point of view. While this definitely resulted in a slower, different style, I feel this approach worked well to contrast my polished 'marketing style' release trailer with a more grounded, accessible introduction to my platform.

Whilst I didn't send out a formal survey regarding this video I did show it to a number of family members and friends who all gave positive feedback on the clarity of my explanations and the engaging nature of my delivery. I feel as if I have demonstrated the features of my platform well and I am phenomenally pleased and proud of the final result. 

The past nine months developing NeuroNote have been an incredible journey and I have learned so much about research techniques, software development, project management and myself. I am excited to see where this project goes in the future, and I am grateful for the opportunity to have worked on something so meaningful and impactful to many of my peers.